\documentclass{beamer}
\beamertemplatenavigationsymbolsempty
\usecolortheme{beaver}
\setbeamertemplate{blocks}[rounded=true, shadow=true]
\setbeamertemplate{footline}[page number]
\usefonttheme{professionalfonts}

\title{An Exploration of Softmax Alternatives Belonging to the Spherical Loss Family}
\author{Alexandre de Brebisson, Pascal Vincent}
\institute{Talk by: Fedor Sobolevsky}
\date{}

\begin{document}

\begin{frame}
\titlepage
\end{frame}

\begin{frame}{Motivation}
\begin{itemize}
    \item Softmax is standard for multi-class classification, but this choice is somewhat arbitrary as its superiority is not theoretically justified.
    \item The spherical loss family allows efficient gradient updates in $O(d^2)$ instead of $O(d \times D)$.
    \item In their paper, Alexandre de Brebisson and Pascal Vincent explore spherical losses as alternatives to log-softmax for classification tasks.
\end{itemize}
\end{frame}

\begin{frame}{Baseline: Log-Softmax}
\[
L(\boldsymbol{o}, c) = -\log\frac{e^{o_c}}{\sum_{k=1}^D e^{o_k}} = -o_c + \log \sum_{k=1}^D e^{o_k}
\]
\begin{itemize}
    \item Standard in classification, but computationally expensive for large $D$: gradient updates require $O(D\times d)$.
\end{itemize}
\end{frame}

\begin{frame}{Spherical Loss Family: Definition}
A loss belongs to the spherical family if it depends only on:
\[
s = \sum_i o_i,\quad q = \|\boldsymbol{o}\|^2,\quad o_c \text{ and } y_c \text{ for target class } c:
\]
\[
\mathcal{L} = \mathcal{L}(s, q, o_c, y_c)
\]
\begin{itemize}
    \item Enables efficient gradient computation\footnote{\textit{Vincent, Pascal, de Brebisson, Alexandre, and Bouthillier, Xavier}. Efficient exact gradient update for training deep networks with very large sparse targets. NIPS, 2015.} in $O(d^2)$.
    \item A classic example is MSE. 
    \end{itemize}
\end{frame}

\begin{frame}{Spherical Upper Bounds of Log-Softmax}
\begin{itemize}
    \item Derived from Bouchard (2007)\footnote{Bouchard, Guillaume. Efficient bounds for the softmax function and applications to approximate inference in hybrid models (2007).} upper bound of log-sum-exp.
    \item Form:
    \[
L \leq \left(-\frac{(D-2)^2}{16D}\frac{1}{\lambda(\xi)}-\frac{D}{2}\xi-D\lambda(\xi)\xi^2+\right.\] \[
\left.+D\log(1+e^{\xi})+
\frac{1}{D}s+\left(q-\frac{s^2}{D}\right)\lambda(\xi)-o_c\right),
\]
\[
\lambda(\xi)=\frac{1}{2\xi}\left(\frac{1}{1+e^{-\xi}}-\frac{1}{2}\right).
\]
    \item Disadvantages:
    \begin{itemize}
        \item Bound minimization did not improve performance.
        \item Hyperparameter $\xi$ sensitive and hard to tune.
        \item Failed to outperform even initial model in language tasks.
    \end{itemize}
\end{itemize}
\end{frame}


\begin{frame}{Spherical Softmax Loss}
\[
f_{\text{sph\_soft}}(\boldsymbol{o})_k = \frac{o_k^2 + \varepsilon}{\sum_i (o_i^2 + \varepsilon)}
\]
\[
L_{\text{log\_sph\_soft}} = -\log f_{\text{sph\_soft}}(\boldsymbol{o})_c
\]
\begin{itemize}
    \item Invariant to scaling of $\boldsymbol{o}$.
    \item Requires careful tuning of $\varepsilon$ for numerical stability.
    \item Even function (ignores sign of $o_k$).
\end{itemize}
\end{frame}

\begin{frame}{Taylor Softmax Loss}
Based on Taylor expansion of $\exp(x) \approx 1 + x + \frac{1}{2}x^2$:
\[
f_{\text{tay\_soft}}(\boldsymbol{o})_k = \frac{1 + o_k + \frac{1}{2}o_k^2}{\sum_i (1 + o_i + \frac{1}{2}o_i^2)}
\]
\[
L_{\text{log\_tay\_soft}} = -\log f_{\text{tay\_soft}}(\boldsymbol{o})_c
\]
\begin{itemize}
    \item No extra hyperparameter needed.
    \item Slightly asymmetric around zero.
    \item Numerically stable.
\end{itemize}
\end{frame}

\begin{frame}{Results: MNIST and CIFAR-10}
\begin{table}[h]
\centering
\begin{tabular}{|l|c|c|}
\hline
\textbf{Loss} & \textbf{MNIST Error} & \textbf{CIFAR-10 Error} \\
\hline
Log-Softmax & 0.812\% & 8.52\% \\
Log-Taylor Softmax & \textbf{0.785\%} & \textbf{8.07\%} \\
Log-Spherical Softmax & 0.828\% & 8.37\% \\
\hline
\end{tabular}
\end{table}
\begin{itemize}
    \item Taylor softmax outperforms log-softmax on low-dimensional outputs.
    \item Spherical losses train faster and achieve lower error.
\end{itemize}
\end{frame}

\begin{frame}{Results: CIFAR-100 and Language Modeling}
\begin{table}[h]
\centering
\begin{tabular}{|l|c|c|}
\hline
\textbf{Loss} & \textbf{CIFAR-100 Error} & \textbf{PennTree Perplexity} \\
\hline
Log-Softmax & \textbf{32.4\%} & \textbf{126.7} \\
Log-Taylor Softmax & 33.1\% & 147.2 \\
Log-Spherical Softmax & 33.1\% & 149.2 \\
\hline
\end{tabular}
\end{table}
\begin{itemize}
    \item Log-softmax performs better as output dimension increases.
    \item Exponential in softmax may better handle high-dimensional competition.
    \item Spherical losses struggle in language tasks with large vocabularies.
\end{itemize}
\end{frame}

\begin{frame}{Key Takeaways}
\begin{itemize}
    \item Spherical losses (especially Taylor softmax) can outperform log-softmax on small-output tasks (e.g., MNIST, CIFAR-10).
    \item For high-dimensional outputs (CIFAR-100, language modeling), log-softmax remains superior.
    \item Spherical losses enable efficient training but may lack the discriminative power of softmax in complex scenarios.
    \item Spherical alternatives are worth considering for specific applications with low output dimension or where efficiency is key.
\end{itemize}
\end{frame}

\end{document}